\item Un taladro eléctrico, con una broca de acero de masa $m = 27.0 \text{ g}$ y $0.635 \text{ cm}$ de diámetro, se usa para taladrar un bloque cúbico de acero de masa $M = 240 \text{ g}$. Suponga que el acero tiene las mismas propiedades que el hierro. Puede modelar el proceso de corte como un punto sobre la circunferencia de la broca. Este punto se mueve en una hélice con rapidez tangencial constante de $40.0 \text{ m/s}$ y ejerce una fuerza de magnitud constante de $3.20 \text{ N}$ sobre el bloque. Como se muestra en la figura P7.8, un surco en la broca lleva las virutas a la parte superior del bloque, donde forman una pila alrededor del orificio. El taladro gira y actúa sobre el bloque durante un intervalo de tiempo de $15.0 \text{ s}$. Suponga que este intervalo es lo suficientemente largo para conducción dentro del acero para llevarlo por completo a una temperatura uniforme. Además, suponga que los objetos de acero pierden una cantidad despreciable de energía por conducción, convección y radiación a sus alrededores. 
(a) Suponga que la broca está afilada y corta tres cuartos del camino a través del bloque durante 15.0 s. Encuentre el cambio de temperatura de toda la cantidad de acero. 
(b) \textbf{¿Qué pasaría si?} Ahora suponga que la broca está mellada y solo corta un octavo del camino a través del bloque durante 15.0 s. Identifique el cambio de temperatura de toda la cantidad de acero en este caso. 
(c) ¿Qué fragmentos de información, si los hay, no son necesarios para la solución? Explique.
